\par\Needspace{6\baselineskip}\medskip
\hypertarget{native-118-article088}{}\pdfbookmark[2]{88.}{article-088}
\begin{center}88.\end{center}
Die Darstellung der \textit{zweiten} Methode gibt uns Gelegenheit, sehr viele neue und elegante Beziehungen herzuleiten. Da diese Beziehungen bei den verschiedenen Arten von Kegelschnitten verschiedene Formen annehmen, müssen wir jede Art für sich behandeln. Wir beginnen mit der ELLIPSE.

\SourceLeaf{119}Den beiden Orten seien die wahren Anomalien $v$ und $v'$ zugeordnet, wobei $v$ zum zeitlich früheren Ort gehöre. Die zugehörigen exzentrischen Anomalien seien $E$ und $E'$, die zugehörigen Radiusvektoren $r$ und $r'$. Ferner sei $p$ der Halbparameter, $e=\sin\varphi$ die Exzentrizität, $a$ die große Halbachse und $t$ die Zeit, in der die Bewegung vom ersten zum zweiten Ort erfolgt. Schließlich setzen wir $v'-v=2f$, $v'+v=2F$, $E'-E=2g$, $E'+E=2G$ und $a\cos\varphi=\dfrac{p}{\cos\varphi}=b$. Durch Kombination der Formeln V und VI aus Art. 8 erhält man nun leicht die folgenden Gleichungen.
\[
\begin{aligned}
\mathrm{1)}\qquad&b\sin g=\sin f\mathbin{.}\BareRoot{rr'}\\[4pt]
\mathrm{2)}\qquad&b\sin G=\sin F\mathbin{.}\BareRoot{rr'}
\end{aligned}
\]
\[
\begin{aligned}
p\cos g={}&\{\cos\tfrac12v\cos\tfrac12v'\mathbin{.}(1+e)\\[3pt]
&\quad+\sin\tfrac12v\sin\tfrac12v'\mathbin{.}(1-e)\}\BareRoot{rr'}
\end{aligned}
\]
oder
\[
\mathrm{3)}\qquad p\cos g=(\cos f+e\cos F)\BareRoot{rr'}
\]
und ebenso
\[
\mathrm{4)}\qquad p\cos G=(\cos F+e\cos f)\BareRoot{rr'}.
\]
Aus der Kombination der Gleichungen 3 und 4 folgt weiter
\[
\begin{aligned}
\mathrm{5)}\qquad&\cos f\mathbin{.}\BareRoot{rr'}=(\cos g-e\cos G)a\\[4pt]
\mathrm{6)}\qquad&\cos F\mathbin{.}\BareRoot{rr'}=(\cos G-e\cos g)a.
\end{aligned}
\]
Aus Formel III von Art. 8 erhält man
\[
\begin{aligned}
\mathrm{7)}\qquad r'-r&=2ae\sin g\sin G\\[4pt]
r'+r&=2a-2ae\cos g\cos G\\[4pt]
&=2a\sin g^2+2\cos f\cos g\mathbin{.}\BareRoot{rr'},
\end{aligned}
\]
und daraus
\[
\mathrm{8)}\qquad a=\frac{r+r'-2\cos f\cos g\mathbin{.}\BareRoot{rr'}}{2\sin g^2}\mathpunct{\cdot}
\]
Setzen wir
\[
\mathrm{9)}\qquad\frac{\BareRoot{\dfrac{r'}r}+\BareRoot{\dfrac{r}{r'}}}{2\cos f}=1+2l,
\]
so erhalten wir
\[
\mathrm{10)}\qquad a=\frac{2(l+\sin\tfrac12g^2)\cos f\mathbin{.}\BareRoot{rr'}}{\sin g^2}
\]
sowie $\BareRoot{a}=\pm\dfrac{\BareRoot{(2(l+\sin\tfrac12g^2)\cos f\mathbin{.}\BareRoot{rr'})}}{\sin g}\mathpunct{\raisebox{.45ex}{,}}$ wobei das obere oder untere Vorzeichen zu wählen ist, je nachdem, ob $\sin g$ positiv oder negativ ist. Formel XII aus Art. 8 \SourceLeaf{120}liefert die Gleichung
\[
\begin{aligned}
\frac{kt}{a^{3/2}}&=E'-e\sin E'-E+e\sin E\\[4pt]
&=2g-2e\sin g\cos G\\[4pt]
&=2g-\sin2g+2\cos f\sin g\frac{\BareRoot{rr'}}a\mathpunct{\cdot}
\end{aligned}
\]
Setzt man nun in dieser Gleichung für $a$ den Wert aus Gleichung 10 ein und schreibt der Kürze halber
\[
\mathrm{11)}\qquad\frac{kt}{2^{3/2}\cos f^{3/2}(rr')^{3/4}}=m,
\]
so erhält man nach entsprechender Umformung
\[
\mathrm{12)}\qquad\pm m=(l+\sin\tfrac12g^2)^{1/2}+(l+\sin\tfrac12g^2)^{3/2}\left(\frac{2g-\sin2g}{\sin g^3}\right),
\]
Dabei ist vor $m$ das obere oder untere Vorzeichen zu setzen, je nachdem, ob $\sin g$ positiv oder negativ ist.

Liegt die heliozentrische Bewegung zwischen $180^0$ und $360^0$ oder ist allgemeiner $\cos f$ negativ, so wäre die nach Formel 11 bestimmte Größe $m$ imaginär und $l$ negativ. Um dies zu vermeiden, verwenden wir in diesem Fall anstelle der Gleichungen 9 und 11 die folgenden Gleichungen.
\[
\begin{aligned}
\mathrm{9^*)}\qquad&\frac{\BareRoot{\dfrac{r'}r}+\BareRoot{\dfrac{r}{r'}}}{2\cos f}=1-2L\\[8pt]
\mathrm{11^*)}\qquad&\frac{kt}{2^{3/2}(-\cos f)^{3/2}(rr')^{3/4}}=M,
\end{aligned}
\]
Daraus erhalten wir anstelle der Gleichungen 10 und 12 die folgenden Gleichungen.
\[
\begin{aligned}
\mathrm{10^*)}\qquad a&=\frac{-2(L-\sin\tfrac12g^2)\cos f\mathbin{.}\BareRoot{rr'}}{\sin g^2}\\[8pt]
\mathrm{12^*)}\qquad\pm M&=-(L-\sin\tfrac12g^2)^{1/2}+(L-\sin\tfrac12g^2)^{3/2}\left(\frac{2g-\sin2g}{\sin g^3}\right),
\end{aligned}
\]
Das Vorzeichen ist hier auf dieselbe Weise zu wählen wie zuvor.
